\SourcePage{235}{240}
\section{Photoelectric effect. Nonrelativistic case}

In \S\S~49--52 we considered radiative transitions (emission or absorption of a photon) between atomic levels belonging to the discrete spectrum. The photoelectric effect differs from such a process of%
\SourcePage{236}{241}
photon absorption only in that the final state belongs to the continuous spectrum.

For a hydrogen atom, or a hydrogen-like ion (with nuclear charge $Z\ll137$), the photoelectric cross section can be calculated analytically in closed form.

In the initial state there is an electron in a discrete level $\varepsilon_i=-I$ ($I$ is the ionization potential of the atom) and a photon of definite momentum $\mathbf k$. In the final state the electron has momentum $\mathbf p$ (and energy $\varepsilon_f=\varepsilon$). Since $\mathbf p$ ranges over a continuous set of values, the photoelectric cross section is
\begin{equation}
d\sigma=2\pi|V_{fi}|^2\delta(-I+\omega-\varepsilon)\frac{d^3p}{(2\pi)^3}
\tag{56.1}
\end{equation}
(cf. (44.3)), where the wave function of the final electron state is assumed to be normalized to one particle in a volume $V=1$. The photon wave function is normalized in the same way as before; in passing from the probability $dw$ to the cross section $d\sigma$, it must therefore be divided by the photon flux density (equal to $c/V=c$), but in relativistic units this does not alter the form of (56.1).

As in (45.2), choose the three-dimensionally transverse gauge for the photon. Then
\[
V_{fi}=-e\mathbf A\mathbf j_{fi}
=-e\sqrt{4\pi}\frac1{\sqrt{2\omega}}M_{fi},
\]
where
\begin{equation}
M_{fi}=\int\psi'^*(\boldsymbol\alpha\mathbf e)e^{i\mathbf k\mathbf r}\psi\,d^3x
\tag{56.2}
\end{equation}
($\psi\equiv\psi_i$ and $\psi'\equiv\psi_f$ are the initial and final electron wave functions). Replacing $d^3p\to\mathbf p^2d|\mathbf p|\,do=\varepsilon|\mathbf p|\,d\varepsilon\,do$ in (56.1) and integrating the $\delta$-function with respect to $d\varepsilon$, we rewrite this formula as
\begin{equation}
d\sigma=e^2\frac{\varepsilon|\mathbf p|}{2\pi\omega}|M_{fi}|^2do.
\tag{56.3}
\end{equation}

We shall carry out the calculation in two cases, distinguished by the photon energy: for $\omega\gg I$ and for $\omega\ll m$. Since $I\sim me^4Z^2\ll m$, these two regions overlap in part (when $I\ll\omega\ll m$), so their examination gives, in essence, a complete description of the photoelectric effect.

We begin with the case
\begin{equation}
\omega\ll m.
\tag{56.4}
\end{equation}
Here the electron velocity is small in both the initial and final states, and the problem is therefore wholly nonrelativistic as regards the electron. Accordingly, in (56.2) replace $\boldsymbol\alpha$ by the nonrelativistic velocity operator $\hat{\mathbf v}=-i\nabla/m$%
\SourcePage{237}{242}
(cf. \S~45). We may also pass to the dipole approximation, putting $e^{i\mathbf k\mathbf r}\approx1$, that is, neglecting the photon momentum in comparison with the electron momentum. We then have
\begin{equation}
d\sigma=e^2\frac{m|\mathbf p|}{2\pi\omega}
|\mathbf e\mathbf v_{fi}|^2do,\qquad
\mathbf v_{fi}=-\frac im\int\psi'^*\nabla\psi\cdot d^3x.
\tag{56.5}
\end{equation}

Consider the photoelectric effect from the ground level of a hydrogen atom (or a hydrogen-like ion). Then
\begin{equation}
\psi=\frac{(Ze^2m)^{3/2}}{\sqrt\pi}e^{-Ze^2mr}
\tag{56.6}
\end{equation}
(in ordinary units $me^2\to1/a_0$, where $a_0=\hbar^2/me^2$ is the Bohr radius).

For $\psi'$, on the other hand, we must use the wave function whose asymptotic form contains a plane wave ($e^{i\mathbf p\mathbf r}$) together with an incoming spherical wave (see III, \S~136, where this function was denoted by $\psi_{\mathbf p}^{(-)}$). By the selection rule for $l$, a transition from an $s$-state is possible only to a $p$-state (in the dipole case). Hence, in the expansion\footnote{Below in this section, $p$ denotes $|\mathbf p|$.}
\begin{equation}
\psi_{\mathbf p}^{(-)}
=\frac1{2p}\sum_{l=0}^{\infty}i^l(2l+1)e^{-i\delta_l}
R_{pl}(r)P_l(\mathbf n\mathbf n_1)
\tag{56.7}
\end{equation}
($\mathbf n=\mathbf p/p$, $\mathbf n_1=\mathbf r/r$), it is sufficient to retain only the term with $l=1$. Dropping inessential phase factors, we obtain
\begin{equation}
\psi'=\frac3{2p}(\mathbf n\mathbf n_1)R_{p1}(r).
\tag{56.8}
\end{equation}
With the functions $\psi$ and $\psi'$ from (56.6) and (56.8), we have
\[
\begin{aligned}
\mathbf e\mathbf v_{fi}
&=\frac{3(Ze^2m)^{5/2}}{2\sqrt\pi\,mp}
\iint(\mathbf n\mathbf n_1)(\mathbf n_1\mathbf e)
e^{-Ze^2mr}R_{p1}(r)\,do_1\cdot r^2dr={}\\
&=\frac{\sqrt{2\pi}(Ze^2m)^{5/2}}{pm}
\int_0^\infty r^2e^{-Ze^2mr}R_{p1}(r)\,dr.
\end{aligned}
\]
According to (36.18) and (36.24) (see III), the radial function (in the units used here) is
\[
R_{pl}=\frac{\sqrt{8\pi}\,Ze^2m}{3}
\sqrt{\frac{1+\nu^2}{\nu(1-e^{-2\pi\nu})}}\,
pr\,e^{-ipr}F(2+i\nu,4,2ipr),
\]
where
\begin{equation}
\nu=\frac{Ze^2m}{p}\left(=\frac{Ze^2}{\hbar v}\right).
\tag{56.9}
\end{equation}
\SourcePage{238}{243}
The integral required here is evaluated with the aid of the formula
\[
\int_0^\infty e^{-\lambda z}z^{\gamma-1}F(\alpha,\gamma,kz)\,dz
=\Gamma(\gamma)\lambda^{\alpha-\gamma}(\lambda-k)^{-\alpha}
\]
(see III, (f. 3)). Noting also that
\[
\left(\frac{\nu+i}{\nu-i}\right)^{i\nu}
=e^{-2\nu\operatorname{arccot}\nu},
\]
we obtain
\[
\mathbf e\mathbf v_{fi}
=\frac{2^{7/2}\pi\nu^3(\mathbf n\mathbf e)}
{\sqrt{pm}(1+\nu^2)^{3/2}}\,
\frac{e^{-2\nu\operatorname{arccot}\nu}}
{\sqrt{1-e^{-2\pi\nu}}}.
\]

The ionization energy from the ground level of a hydrogen atom (or hydrogen-like ion) is $I=Z^2e^4m/2$. Therefore
\begin{equation}
\omega=\frac{p^2}{2m}+I=\frac{p^2}{2m}(1+\nu^2).
\tag{56.10}
\end{equation}
Taking this relation into account, the final expression for the photoelectric cross section with emission of an electron into the solid-angle element $do$ is
\begin{equation}
d\sigma=2^7\pi\alpha a^2
\left(\frac I{\hbar\omega}\right)^{\!4}
\frac{e^{-4\nu\operatorname{arccot}\nu}}
{1-e^{-2\pi\nu}}(\mathbf n\mathbf e)^2do,
\tag{56.11}
\end{equation}
where $a=\hbar^2/(mZe^2)=a_0/Z$ (here and below, ordinary units are used). Notice that the angular distribution of the photoelectrons is determined by the factor $(\mathbf n\mathbf e)^2$. It is greatest in directions parallel to the polarization of the incident photons, and vanishes in directions perpendicular to $\mathbf e$, including the direction of incidence. For unpolarized photons, (56.11) must be averaged over the directions of $\mathbf e$, which amounts to the replacement
\[
(\mathbf n\mathbf e)^2\to\frac12[\mathbf n_0\mathbf n]^2,
\]
where $\mathbf n_0=\mathbf k/k$ (see (45.46)).

Integration of (56.11) over the angles gives the total photoelectric cross section:
\begin{equation}
\sigma=\frac{2^9\pi^2}{3}\alpha a^2
\left(\frac I{\hbar\omega}\right)^{\!4}
\frac{e^{-4\nu\operatorname{arccot}\nu}}
{1-e^{-2\pi\nu}}
\tag{56.12}
\end{equation}
(\emph{M. Stobbe}, 1930).

The limiting value of $\sigma$ as $\hbar\omega\to I$ (that is, $\nu\to\infty$) is
\begin{equation}
\sigma=\frac{2^9\pi^2}{3e^4}\alpha a^2
=\frac{2^9\pi^2}{3e^4}\frac{\alpha a_0^2}{Z^2}
=0{.}23\frac{a_0^2}{Z^2}
\tag{56.13}
\end{equation}
\SourcePage{239}{244}
(in the denominator, $e=2{.}71\ldots$!). As is required for a reaction producing charged particles (see III, \S~147), the photoelectric cross section tends to a constant limit near its threshold.

The case $\hbar\omega\gg I$ (while still $\hbar\omega\ll mc^2$) corresponds to the Born approximation ($\nu=Ze^2/(\hbar v)\ll1$). Formula (56.12) then becomes
\begin{equation}
\sigma=\frac{2^8\pi}{3}\alpha a_0^2Z^5
\left(\frac{I_0}{\hbar\omega}\right)^{\!7/2}
\tag{56.14}
\end{equation}
($I_0=e^4m/(2\hbar^2)$ is the ionization energy of the hydrogen atom).

The inverse of the photoelectric effect is radiative recombination of an electron with a stationary ion. The cross section for this process ($\sigma_{\text{rec}}$) can be found from the photoelectric cross section ($\sigma_{\text{ph}}$) by means of the principle of detailed balance (III, \S~144). According to this principle, the cross sections of the processes $i\to f$ and $f\to i$ (with two particles in each of the states $i$ and $f$) are related by
\[
g_i p_i^2\sigma_{i\to f}=g_f p_f^2\sigma_{f\to i},
\]
where $p_i$, $p_f$ are the relative-motion momenta of the particles, and $g_i$, $g_f$ are the spin statistical weights of states $i$ and $f$. Taking account also of the fact that $g=2$ for a photon (two polarization directions), while the statistical weight of a free electron and ion is equal to that of the ground state of the hydrogen atom, we obtain for this state
\begin{equation}
\sigma_{\text{rec}}=\sigma_{\text{ph}}\frac{2k^2}{p^2}
\tag{56.15}
\end{equation}
($\mathbf p=m\mathbf v$ is the momentum of the incident electron, and $\mathbf k$ is the momentum of the emitted photon).

\ProblemHeading

1. Derive (56.14) by applying the Born approximation directly in the nonrelativistic case.

\SolutionHeading In the Born approximation, $\psi'$ in (56.5) must simply be taken as the plane wave $\psi'=e^{i\mathbf p\mathbf r}$, while $\psi$ remains the function (56.6). Then
\[
\mathbf v_{fi}=\mathbf v_{if}
=\frac1m\int\psi\hat{\mathbf p}\psi'\,d^3x
=\frac{\mathbf p}{m}\frac{(Ze^2m)^{3/2}}{\sqrt\pi}
(e^{-Ze^2mr})_{\mathbf p}.
\]
The Fourier component is given by (57.66), and hence
\[
\mathbf v_{fi}\approx8\sqrt\pi\,p^{-3}m^{3/2}(Ze^2)^{5/2}\mathbf n.
\]
Substitution in (56.5), followed by integration over $do$, gives (56.14) (to the required accuracy, $p^2/(2m)\approx\omega$ here).

2. Find the total cross section for radiative recombination of a fast nonrelativistic electron ($I\ll mv^2\ll mc^2$) with a nucleus (of charge $Z\ll137$).

\SourcePage{240}{245}
\SolutionHeading The capture cross section into the $K$ shell (principal quantum number $n=1$) follows by substituting (56.14) into (56.15):
\[
\sigma_1^{\text{rec}}
=\frac{2^7\pi}{3}Z^5\alpha^3a_0^2
\left(\frac{I_0}{\varepsilon}\right)^{\!5/2}
\]
($\varepsilon=mv^2/2$ is the energy of the incident electron; $\hbar\omega\approx\varepsilon$). Of the other states of the atom being formed, only the $s$-states are important: in calculating the matrix element in the Born approximation, the values of the bound-state wave function at small $r$ are relevant (as will be seen from the calculations in \S~57), and for $l>0$ these values are small compared with those of the functions having $l=0$; it is sufficient here to retain the first two terms in the expansion of $\psi$ in powers of $r$. For states with $l=0$ and arbitrary $n$, these terms are
\[
\psi=\frac1{\sqrt{\pi a^3}n^{3/2}}
\left(1-\frac ra\right),
\]
that is, they contain $n$ only through the overall factor $n^{-3/2}$ (the displayed expression follows by expanding (36.13) (see III)). Consequently, the total recombination cross section is
\[
\sigma^{\text{rec}}
=\sum_{n=1}^{\infty}\sigma_n^{\text{rec}}
=\sigma_1^{\text{rec}}\sum_{n=1}^{\infty}\frac1{n^3}
=\zeta(3)\sigma_1^{\text{rec}}
\]
(the value of the zeta function is $\zeta(3)=1{.}202$).
